\chapter{Dane i metodyka badań}
\label{chap:methodology}
\section{Źródła danych i ich zgodność}
\writingnote{Opisz dane publiczne, syntetyczne i fizyczne: pochodzenie, licencje, jednostki, cechy, etykiety i ograniczenia. Zbiory mogą służyć odrębnym eksperymentom, jeżeli ich semantyka jest niezgodna.}
\section{Przygotowanie danych i cech}
\writingnote{Zdefiniuj okna czasowe, brakujące pomiary, skalowanie i relacje sensor--aktuator. Osobno traktuj ruch obserwowany przez bramę i wartości raportowane przez urządzenie.}
\section{Podział danych i zapobieganie wyciekowi informacji}
\writingnote{Rozdziel niezależne przebiegi lub okresy. Dopasuj preprocessing wyłącznie do treningu; strojenie wykonuj na walidacji. Końcowych danych fizycznych do testu nie używaj do strojenia.}
\section{Metody porównawcze i hipotezy}
\writingnote{Wybierz reguły odniesienia i 2--3 algorytmy. Zapisz hipotezy przed interpretacją wyników, zakres strojenia oraz wspólny budżet eksperymentalny.}
\section{Metryki skuteczności i wydajności}
\writingnote{Zdefiniuj precision, recall, F1, FPR, fałszywe alarmy na godzinę i czas wykrycia. Ustal jednostkę oceny: próbka lub zdarzenie, grupowanie alarmów i traktowanie niewykrytych zdarzeń.}
\writingnote{Dla opóźnienia inferencji, CPU i RAM podaj sposób pomiaru, rozgrzewkę, liczbę powtórzeń oraz charakterystykę sprzętu. Rozdziel opóźnienie samego modelu od całego przetwarzania.}
\section{Plan eksperymentów i odtwarzalność}
\begin{table}[htbp]
\centering\small
\caption{Proponowane grupy eksperymentów. Są to plany, nie wyniki.}
\label{tab:experiment-plan}
\begin{tabularx}{\textwidth}{@{}lXl@{}}
\toprule
Grupa & Porównanie & Pytania \\
\midrule
E1 & Reguły i wybrane metody ML przy wspólnym podziale danych & RQ1, RQ5 \\
E2 & Przeniesienie modelu na zgodne dane fizycznego prototypu & RQ2 \\
E3 & Detekcja dla różnych węzłów i zestawów możliwości & RQ3 \\
E4 & Zmiana liczby węzłów i intensywności wiadomości & RQ4 \\
E5 & Opóźnienie, pamięć i CPU na Raspberry Pi & RQ1, RQ4 \\
\bottomrule
\end{tabularx}
\end{table}

\writingnote{Dla każdego przebiegu zachowaj identyfikator, konfigurację, seed, commit kodu, wersję danych i modelu. Dobierz liczby węzłów i powtórzenia po pomiarach możliwości sprzętu.}
